\kbvexercise{ಅಭ್ಯಾಸ ಪತ್ರಿಕೆ ೬}

\vspace{2.5pt plus 3pt minus 0.5pt}\noindent\hspace*{1.2em}\kbvnum{(1)}\hspace{0.55em}ಸಂಧಿ ಎಂದರೇನು?\par

\vspace{2.5pt plus 3pt minus 0.5pt}\noindent\hspace*{1.2em}\kbvnum{(2)}\hspace{0.55em}ಕನ್ನಡ ಸಂಧಿಗಳಲ್ಲಿ ಎಷ್ಟು ಪ್ರಕಾರಗಳಿವೆ?\par

\vspace{2.5pt plus 3pt minus 0.5pt}\noindent\hspace*{1.2em}\kbvnum{(3)}\hspace{0.55em}ಕನ್ನಡ ಸಂಧಿಗಳಲ್ಲಿ ಸ್ವರ ಸಂಧಿಗಳು ಯಾವುವು?\par

\vspace{2.5pt plus 3pt minus 0.5pt}\noindent\hspace*{1.2em}\kbvnum{(4)}\hspace{0.55em}ವ್ಯಂಜನ ಸಂಧಿ ಎಂದರೇನು? ಉದಾಹರಣೆ ಸಹಿತ ತಿಳಿಸಿ.\par

\vspace{2.5pt plus 3pt minus 0.5pt}\noindent\hspace*{1.2em}\kbvnum{(5)}\hspace{0.55em}ಲೋಪ ಸಂಧಿ ಎಂದರೇನು? ಉದಾಹರಣೆ ಸಹಿತ ವಿವರಿಸಿ.\par

\vspace{2.5pt plus 3pt minus 0.5pt}\noindent\hspace*{1.2em}\kbvnum{(6)}\hspace{0.55em}ಆಗಮ ಸಂಧಿ ಎಂದರೇನು? ಉದಾಹರಣೆ ಸಹಿತ ವಿವರಿಸಿ.\par

\vspace{2.5pt plus 3pt minus 0.5pt}\noindent\hspace*{1.2em}\kbvnum{(7)}\hspace{0.55em}ಆದೇಶ ಸಂಧಿ ಎಂದರೇನು? ಉದಾಹರಣೆ ಸಹಿತ ವಿವರಿಸಿ.\par

\vspace{2.5pt plus 3pt minus 0.5pt}\noindent\hspace*{1.2em}\kbvnum{(8)}\hspace{0.55em}ಯ-ಕಾರಾಗಮ ಸಂಧಿಗೆ ಎರಡು ಉದಾಹರಣೆಗಳನ್ನು ಕೊಡಿ.\par

\vspace{2.5pt plus 3pt minus 0.5pt}\noindent\hspace*{1.2em}\kbvnum{(9)}\hspace{0.55em}-- ಇವುಗಳನ್ನು ಕೂಡಿಸಿ ಬರೆದು, ಸಂಧಿಯ ಹೆಸರು ತಿಳಿಸಿ.\par

\begin{kbvderiv}
\vspace{4pt plus 3pt minus 1pt}
\noindent{\renewcommand{\arraystretch}{1.24}\setlength{\tabcolsep}{5pt}
\begin{tabularx}{\textwidth}{X X X}
ಹಸಿರು + ಎಲೆ & ಹಣೆ + ಕಣ್ಣು & ಕೆಂಪು + ತುಟಿ \\
ಈಗ + ಈಗ & ಗಿಡ + ಇದು & ತಂತಿ + ಇಂದ \\
ಮೈ + ತಡವು & ಬಾರೆ + ಅಮ್ಮ & ನನಗೆ + ಇದೆ \\
\end{tabularx}}
\vspace{4pt plus 3pt minus 1pt}
\end{kbvderiv}

\vspace{2.5pt plus 3pt minus 0.5pt}\noindent\hspace*{1.2em}\kbvnum{(10)}\hspace{0.55em}ಈ ಸಂಧಿ-ಪದಗಳನ್ನು ಬಿಡಿಸಿ ಬರೆದು, ಸಂಧಿಯ ಹೆಸರು ತಿಳಿಸಿ.\par

\begin{kbvderiv}
\vspace{4pt plus 3pt minus 1pt}
\noindent{\renewcommand{\arraystretch}{1.24}\setlength{\tabcolsep}{5pt}
\begin{tabularx}{\textwidth}{X X X}
ಊರನ್ನು & ಹುಲ್ಲುಗಾವಲು & ಮನೆಯನ್ನು \\
ದೇಶವನ್ನು & ಆಗಾಗ & ಮಳೆಗಾಲ \\
\end{tabularx}}
\vspace{4pt plus 3pt minus 1pt}
\end{kbvderiv}

\kbvchapter{ಪಾಠ ೭}{ಸಂಸ್ಕೃತ ಸಂಧಿಗಳು}

\indent ಕನ್ನಡ ಭಾಷೆಯಲ್ಲಿ ಬಹುತೇಕ ಸಂಸ್ಕೃತ ಶಬ್ದಗಳು ಸೇರಿಹೋಗಿವೆ. ಇಂತಹ ಸಂಸ್ಕೃತ ಶಬ್ದಗಳಲ್ಲಿ ಎರಡು ಶಬ್ದಗಳು ಸೇರಿ ಸಂಧಿಯಾಗಿರುವ ಪಕ್ಷಗಳು ಸಾಕಷ್ಟಿವೆ. ಇಲ್ಲಿ ಸಂಧಿಕಾರ್ಯವು ಸಂಸ್ಕೃತ ಭಾಷಾಶಾಸ್ತ್ರಕ್ಕೆ ಅನುಗುಣವಾಗಿ ನಡೆಯುವ ಕಾರಣ ಸಂಸ್ಕೃತ ಸಂಧಿಗಳನ್ನು ತಿಳಿಯಬೇಕಾದ್ದು.\par

\indent ಸಂಸ್ಕೃತ ಸಂಧಿಗಳನ್ನು ಸ್ವರ ಸಂಧಿಗಳು ಮತ್ತು ವ್ಯಂಜನ ಸಂಧಿಗಳು ಎಂದು ಎರಡು ಬಗೆಯಾಗಿ ವಿಂಗಡಿಸಬಹುದು. ಇವುಗಳಲ್ಲಿ (1) ಸವರ್ಣ ಸಂಧಿ, (2) ಗುಣ ಸಂಧಿ, (3) ವೃದ್ಧಿ ಸಂಧಿ, ಮತ್ತು (4) ಯಣ್‌\kbvcontinues
